\scanpage{110}

\diagram{\begin{tikzcd}[ampersand replacement=\&,column sep=small,row sep=small]
\mathbb Q'(i)\arrow[r,dash]\arrow[d,dash] \& L\arrow[d, dash]\arrow[dr,dash] \& \deleted{\unsure{\mathbb Q'(i)}}\\
\mathbb Q'\arrow[r,dash] \& K\arrow[dr,dash]\arrow[dd,dash] \& L_0\arrow[d,dash] \&[-2ex] \bigg)\ \text{\unsure{galoisienne|galoisiens}}\\
\& \& K_0\\
\& \mathbb Q
\end{tikzcd}}

On veut construire le gpe de \unsure{t.m.} \inserted{$G(q)$} associé
à l'ens $M(q)$ des $\lambda\in L$ tels que $\lambda\bar\lambda=q^i$,
et $\lambda$ entiers exc. $p$. On a vu que

\[
M_0(q)/\mathrm{Torsion}\unsure{\subset|-|\simeq}(\mathbb Z^P)'\tag*{\unsure{(\ensuremath{*})|(\ensuremath{\pi})}}
\]

% slanted
\margin{N.B. $M_0(q)\inserted{=M_0(p)}$ \unsure{(ens des|= ens. des} $u\in L$ qui sont des unités exc. $p$ et satisfont $u\bar u=1$.}
$P$ l'ens des idéaux premiers de $L$ sur $p$.
On \unsure{peut supposer|doit supposer|l'a supposé} que les idéaux $\unsure{\mathfrak p|p|\text{pre}}$ de $K$ se \uline{décomposent} dans $L$.
Choisissons un des \deleted{\unsure{ces}} idéaux \unsure{$\mathfrak p$ au-dessus de $p$}, ce qui donne
une forme linéaire can. sur $M$

\[
\lambda\unsure{\longmapsto|\longrightarrow}v_{\mathfrak p}(\lambda),
\]

qui engendre d'ailleurs $\check M$ mod torsion sous \deleted{\unsure{$G$}} le
gpe de Galois $\Gamma$ de $L$, d'où

\[
{\mathbb G_m}_L\xrightarrow{v_L}G_L
\]

Soit $\Gamma_{\mathfrak p}$ le \unsure{grpe|gpe} de décomposition pour $\mathfrak p$, \deleted{qui est
\unsure{aussi}} et soit \deleted{$K_0$ le corps} $L_0$ le corps des
invariants de $\Gamma_{\mathfrak p}$, $K_0=L_0\cap K$. Alors
$L_0$ est le plus petit corps de rationalité de $v$,
qui s'écrit \unsure{donc|dans}

\[
v:{\mathbb G_m}_{L_0}\longrightarrow G_{L_0}
\]

D'autre part, si $\mathfrak p_0$ est induit par $\mathfrak p$ sur $L_0$,
alors \deleted{\unsure{tout} l'anneau de valuation complété} \inserted{l'anneau local complété}
de $\mathfrak p_0$ est \deleted{égal à celui de} isom : à celui de
$\mathbb Z_{\ell}$. Notons que le couple $(L_0,\mathfrak p_0)$ est défini
à isom unique près par la donnée de la \uline{classe}
de $L$, et peut se définir \deleted{\unsure{comme}} \inserted{via} \unsure{le|la} hensélis\unsure{é|e}
relatif de $\mathbb Z_{(p\mathbb Z)}$ dans $L$ \ldots. Noter que $L_0$ est
purement imaginaire, grâce à l'hypothèse
que $\tau\mathfrak p\ne\mathfrak p$.

% archive sheet number: 109
